% Section 6 — Evidence for data-conditioned programs
\section{Evidence today: programs written for each dataset}
\label{sec:evidence}

\Cref{sec:frontier} measured route~1, one general program per model. This section turns to route~2:
programs written for one dataset at a time. Everything in it is a pilot or a reference measurement. Each result carries its track
(\cref{tab:tracks}) and an evidence label: \est{} established in the repository with a committed,
reproducible record; \pre{} preliminary, meaning a small panel, an unmerged branch, an estimate or a
run on a non-qualifying host; \pend{} not yet measured. \Cref{app:claims} traces every number to its
file.

\subsection{What the program space contains}
\label{sec:space}

Two reference measurements put those numbers in perspective. Neither is a fair comparison, and we
use neither as one.

On the 64 families of \cref{sec:u64}, the best program that the engine's offline search had previously found for
each family --- selected on the same inner validation, after a much larger and unequal search budget
--- has a geometric-mean RMSE 2.2 to 3.0 times lower than the frontier entrants': the entrants'
errors are 2.29 (Opus), 2.21 (Sol) and 2.97 (Astra) times the incumbent's, and they beat it on 21, 11
and 3 of the 64 families respectively \pre. This is a conditional reference, not a contest: the
incumbent saw far more evaluations and was selected on the same folds it is scored on. What it shows
is how much better the best known program for each dataset is than any single general algorithm,
however carefully reasoned.

The public Golden45 ladder makes the same point from the other side \est. When the engine's best
program per dataset is chosen by looking at test scores --- the minimum over 63,108 test-scored
evaluations on 42 datasets --- it ranks first of 304 methods at Elo 2259.8. A single fixed-seed run of
the engine, which never sees test labels, ranks 31st at 1687.0 with 42 of 45 datasets covered,
below tabular foundation models such as TabFM
(1934.9) and TabPFN v3 (1909.5) \citep{hollmann2025tabpfn}. The first number is an oracle envelope
and must never be read as a result; the gap between the two is the clearest statement we have of
the problem InnovationZero addresses. \emph{Strong programs for these datasets exist in the space.
Finding them without looking at the answer is the whole problem.}

\subsection{Tracks U and S on identical cuts}
\label{sec:pilot}

The repository's committed same-cut comparison ranked the engine third of 15 (Elo 1556.6), but on
ten favourably selected datasets, with only the top five methods deepened and half of the finalists
covered; it contained no frontier entrant. We therefore ran a complete one \pre: the three
merged entrants, the engine's track-S proposer at $k=256$ (32 candidates in each of 8 rounds,
re-run on every cut's training rows), its ten single-head baselines and four external libraries at
their defaults (XGBoost 2.1.4, LightGBM 4.6.0, CatBoost 1.2.10 and scikit-learn's random forest
1.7.2) \citep{chen2016xgboost,ke2017lightgbm,prokhorenkova2018catboost,breiman2001rf}. All eighteen
methods were fitted on the same training rows and scored on the same held-out rows of five cuts of
the repository's ten test datasets, by one evaluator; every method was scored on every cut, so
coverage is complete (\cref{app:pilot} gives the protocol). The datasets are a favourable selection
for the engine: they were chosen, historically, because it ranked first on them.

\begin{figure}[t]
\centering
% Pilot same-cut ladder. Source: evidence/pilot-2026-10-22/elo-deep.csv
% (jope-arena elo, regularized Bradley--Terry; scores.csv SHA-256 86fec779...).
\begin{tikzpicture}
\begin{axis}[iz,
  width=0.74\linewidth, height=0.66\linewidth,
  xbar=0pt, bar width=6.2pt, bar shift=0pt,
  xmin=1300, xmax=1800, xtick={1300,1400,1500,1600,1700,1800},
  xmajorgrids, xlabel={Elo (pilot, 18 methods, 10 datasets $\times$ 5 cuts)},
  ytick={1,...,18}, y dir=reverse, ymin=0.4, ymax=18.6,
  yticklabels={{Opus5p5\_Extra},{Stack, engine baseline},{GPT6p1\_Sol\_Max},{Engine, $k=256$},{Kernel ridge, engine baseline},{HistGBM2, engine baseline},{Ensemble, engine baseline},{GPT6\_Astra\_Max},{Gaussian process, engine baseline},{CatBoost 1.2.10},{Random forest 1.7.2},{HistGBM, engine baseline},{LightGBM 4.6.0},{Elastic net, engine baseline},{Ridge-GCV, engine baseline},{Ridge, engine baseline},{XGBoost 2.1.4},{Kernel smoother, engine baseline}},
  yticklabel style={font=\sffamily\scriptsize\color{izink}},
  axis y line*=left, axis x line*=bottom, clip=false,
  legend style={at={(0.99,0.03)}, anchor=south east, legend columns=1, fill=white, draw=none},
  legend cell align=left]
\addplot+[fill=izwarm, draw=none] coordinates {(1761.4,1) (1649.0,3) (1506.2,8)};
\addplot+[fill=izaccent, draw=none] coordinates {(1630.9,4)};
\addplot+[fill=izgreen!80, draw=none] coordinates {(1482.2,10) (1466.2,11) (1454.1,13) (1326.0,17)};
\addplot+[fill=izrule, draw=none] coordinates {(1691.8,2) (1573.2,5) (1532.3,6) (1522.3,7) (1490.2,9) (1460.2,12) (1383.9,14) (1375.3,15) (1368.8,16) (1326.0,18)};
\legend{Frontier-model entrant (track U), Engine (track S), Library default, Engine single-head baseline}
\foreach \e/\r in {1761/1,1692/2,1649/3,1631/4,1573/5,1532/6,1522/7,1506/8,1490/9,1482/10,1466/11,1460/12,1454/13,1384/14,1375/15,1369/16,1326/17,1326/18}{
  \edef\tmp{\noexpand\node[font=\noexpand\sffamily\noexpand\tiny\noexpand\color{izgrey}, anchor=west] at (axis cs:\e,\r) {\e};}\tmp}
\end{axis}
\end{tikzpicture}

\caption{Tracks U and S on identical cuts: Elo of 18 methods over ten datasets and five cuts
(\code{jope-arena}, regularized Bradley--Terry). The frontier entrants are track U, the engine track
S with $k=256$. Pilot on a 2-vCPU host; non-qualifying \pre.}
\label{fig:ladder}
\end{figure}

\Cref{fig:ladder} gives the result. Opus's regressor ranks first (Elo 1761.4), the engine's own
stacking baseline second (1691.8), Sol third (1649.0) and the engine's searched programs fourth
(1630.9); Astra is eighth (1506.2). The four libraries at their defaults rank tenth to seventeenth.
On error magnitude the picture shifts again: the engine's programs have the lowest geometric-mean
RMSE relative to each dataset's best method among the non-baseline methods (1.157, against 1.171 for
Opus, 1.224 for CatBoost and 1.321 for Sol). Paired over the ten datasets (\cref{tab:pairs}), no
difference between a frontier entrant and the engine is statistically distinguishable; Opus beats
XGBoost and LightGBM on all ten datasets and CatBoost on nine; and Opus and Astra differ clearly
from each other. Per-dataset ranks (\cref{fig:heatmap}) show why aggregates hide so much: most
methods' ranks swing by ten places or more across datasets, while Opus's regressor is the most
consistent, never worse than seventh.

\begin{table}[t]
\caption{Paired comparisons over the ten pilot datasets: wins and losses on pooled out-of-fold RMSE,
two-sided sign-test $p$, and the geometric-mean RMSE ratio (first over second; below 1 favours the
first) \pre.}
\label{tab:pairs}
\centering\sffamily\footnotesize\lining
\begin{tabular}{@{}lrrr@{\qquad}lrrr@{}}
\toprule
Pair & W--L & $p$ & Ratio & Pair & W--L & $p$ & Ratio \\
\midrule
Opus / engine & 6--4 & 0.75 & 1.012 & Engine / XGBoost & 9--1 & 0.021 & 0.765 \\
Sol / engine & 6--4 & 0.75 & 1.142 & Engine / LightGBM & 9--1 & 0.021 & 0.754 \\
Astra / engine & 5--5 & 1.00 & 1.222 & Engine / CatBoost & 8--2 & 0.11 & 0.945 \\
Stack / engine & 5--5 & 1.00 & 0.944 & Opus / XGBoost & 10--0 & 0.002 & 0.774 \\
Opus / Sol & 6--4 & 0.75 & 0.886 & Opus / LightGBM & 10--0 & 0.002 & 0.763 \\
Opus / Astra & 9--1 & 0.021 & 0.829 & Opus / CatBoost & 9--1 & 0.021 & 0.957 \\
Sol / Astra & 7--3 & 0.34 & 0.935 & Sol / CatBoost & 7--3 & 0.34 & 1.080 \\
\bottomrule
\end{tabular}
\end{table}

\begin{figure}[t]
\centering
% Per-dataset rank among 18 methods (1 = lowest pooled out-of-fold RMSE).
% Source: data/pilot-ranks.csv, derived from evidence/pilot-2026-10-22/scores-deep.csv.
\begin{tikzpicture}[x=0.072\linewidth, y=0.46cm]
\def\methods{{"AFM-Lite","Opus","Sol","Astra","Stack","CatBoost","Rand.\ forest","LightGBM","XGBoost"}}
\def\datasets{{"047e62b5","1cb97b1d","21fbfe8b","30ef3eaa","43b4cc3a","693890aa","79e600bc","b5beff4a","c32fa09b","c7d81571"}}
\def\R{{{1,7,6,12,5,10,11,13,15},{3,6,13,17,4,12,2,9,10},{14,1,2,12,7,4,8,5,3},{3,4,2,1,5,10,14,12,15},{9,5,3,8,2,14,16,12,18},{4,7,11,9,3,5,6,12,8},{5,1,12,14,7,4,2,6,13},{8,3,6,12,9,13,16,10,18},{9,2,1,6,4,15,17,14,18},{7,1,3,2,5,12,11,13,18}}}
\foreach \d in {0,...,9}{
  \pgfmathsetmacro{\dn}{\datasets[\d]}
  \node[anchor=east, font=\ttfamily\scriptsize\color{izgrey}] at (-0.15,-\d) {\dn};
  \foreach \m in {0,...,8}{
    \pgfmathtruncatemacro{\r}{\R[\d][\m]}
    \pgfmathsetmacro{\shade}{100-(\r-1)*100/17}
    \fill[izaccent!\shade!white] (\m+0.04,-\d-0.44) rectangle (\m+0.96,-\d+0.44);
    \pgfmathsetmacro{\txt}{\r<7 ? "white" : "izink"}
    \node[font=\sffamily\scriptsize\color{\txt}] at (\m+0.5,-\d) {\r};
  }
}
\foreach \m in {0,...,8}{
  \pgfmathsetmacro{\mn}{\methods[\m]}
  \node[anchor=south west, rotate=35, font=\sffamily\scriptsize] at (\m+0.25,0.55) {\mn};
}
\draw[izwarm, line width=0.8pt] (1,0.5) rectangle (4,-9.5);
\node[font=\sffamily\scriptsize\color{izwarm}, anchor=north] at (2.5,-9.6) {frontier entrants};
\draw[izaccent, line width=0.8pt] (0,0.5) rectangle (1,-9.5);
\node[font=\sffamily\scriptsize\color{izaccent}, anchor=north] at (0.5,-9.6) {engine};
\end{tikzpicture}

\caption{Rank of each method on each pilot dataset among all eighteen (1 = lowest pooled
out-of-fold RMSE; darker is better). The engine and the frontier entrants trade places dataset by
dataset; the default libraries are rarely near the top \pre.}
\label{fig:heatmap}
\end{figure}

We draw three conclusions, none of which depends on the small sample being representative. First,
reasoning-only invention is strong: a program written by a language model that never saw a dataset
ranked first among eighteen, ahead of every library default. Second, on these datasets a
data-conditioned program found with 256 evaluations is about as good as the best frontier entrant,
and better on magnitude than most --- which is a statement about track S, not about track C. Third,
the engine's own stacking head beats its searched programs on error magnitude (ratio 0.944): even
with 256 evaluations, choosing the best program for a small dataset is noisy. That is precisely the
job a learned first proposal would have to do better.

\subsection{Ten datasets, ten algorithms}
\label{sec:code}

Where each frontier entrant wrote one algorithm for every dataset, the engine wrote a different one
for each (\cref{tab:programs}): ridge regression after robust scaling for one, a Gaussian process
with a radial-basis-plus-linear kernel over sigmoid features and pairwise products for another,
gradient-boosted trees after Chebyshev features for a third. Each program was found in 0.33--10.3
seconds on two CPU cores, and the same program was found on both hosts where the search was run. The
components are known methods, documented with their origins in the emitted code \citep{hoerl1970ridge,friedman2001gbm};
what is specific to each dataset is the composition. One program contains a redundant step (robust
scaling applied twice); we note it rather than tidy it.

\begin{table}[t]
\caption{The engine's program for each pilot dataset (screen cut), described by its components;
the full emitted Rust for all ten is in the code supplement \est.}
\label{tab:programs}
\centering\sffamily\footnotesize\lining
\begin{tabularx}{\linewidth}{@{}lrLrr@{}}
\toprule
Dataset & Rows $\times$ features & Program & Found in & Rust \\
\midrule
\code{047e62b5} & 195 $\times$ 12 & robust scaling, standardization, ridge & 0.4 s & 13.7 KB \\
\code{1cb97b1d} & 166 $\times$ 10 & sigmoid features, top-5 pairwise products, Gaussian process (rbf + linear) & 10.3 s & 34.3 KB \\
\code{21fbfe8b} & 1000 $\times$ 21 & winsorizing, Chebyshev degree 6, pairwise products, Gaussian process (rbf + linear) & 1.5 s & 37.9 KB \\
\code{30ef3eaa} & 109 $\times$ 30 & standardization, ranks, polynomial interactions, ridge & 1.1 s & 15.5 KB \\
\code{43b4cc3a} & 482 $\times$ 39 & robust scaling, gradient-boosted trees (depth 6, 89 rounds) & 1.4 s & 16.6 KB \\
\code{693890aa} & 202 $\times$ 1 & target-binned encoding (12 bins), Gaussian process & 9.0 s & 21.9 KB \\
\code{79e600bc} & 180 $\times$ 3 & polynomial interactions, robust scaling, gradient-boosted trees (depth 3) & 0.3 s & 17.9 KB \\
\code{b5beff4a} & 575 $\times$ 15 & ranks, normal scores, ridge with cross-validated penalty & 0.7 s & 24.4 KB \\
\code{c32fa09b} & 1000 $\times$ 9 & ranks, polynomial interactions, bagged histogram boosting & 0.6 s & 24.4 KB \\
\code{c7d81571} & 1000 $\times$ 7 & winsorizing, Chebyshev degree 4, robust scaling, gradient-boosted trees & 0.6 s & 19.0 KB \\
\bottomrule
\end{tabularx}
\end{table}

\paragraph{The emitted code} All ten programs were emitted as standalone Rust and Python, and all
compile: the Rust with \code{rustc --edition=2021 -O -D unsafe-code} and no warnings, the Python
with \code{py\_compile} \est. Neither uses anything beyond its standard library; the Python imports
only \code{csv}, \code{math} and \code{sys}. Two qualifications matter. The C\texttt{++}17 and
JavaScript emitters cover fewer heads: they produced code for the two ridge programs and refused the
other eight, so the honest coverage is 10/10 for Rust and Python and 2/10 for C\texttt{++}17 and
JavaScript (24 of 40 language targets). And each emitted program is a command-line tool that reads
CSV files, so it uses the standard library's file, environment and process modules, which the
challenge forbids its entrants; a fit/predict library form, comparable to the challenge interface,
is part of the protocol. The emitted programs are 13.7--37.9 kilobytes of Rust, the same scale as
the frontier entrants' 25--97 kilobytes, which makes the two code appendices a fair side-by-side.

\subsection{What is not yet shown}

No result in this section is a track-C result. The engine's programs were chosen after executing
256 candidates per cut; the frontier entrants are universal programs, not proposals conditioned on
the data; and no neural proposer has yet committed a first program on held-out data
(\cref{sec:model}). The question in \cref{sec:intro} is therefore still open, and the next section
specifies the experiment that will answer it.
